\chapter{Evaluación de usabilidad} \label{Anexo:CuestionarioUsabilidad}
%\section{Cuestionario de usabilidad}

Para analizar la usabilidad de la aplicación propuesta se utilizará un cuestionario. Dicho cuestionario permitirá valorar diferentes características relacionadas con los requisitos no funcionales del sistema, así como la opinión inicial de los posibke usuarios.

Las dimensiones evaluadas serán las siguientes:

\begin{itemize}
    \item \textbf{Eficiencia}: grado en el que la aplicación permite alcanzar los objetivos propuestos con un bajo coste temporal y computacional.
    \item \textbf{Efectividad}: capacidad de la aplicación para cumplir correctamente los requisitos definidos en el presente trabajo.
    \item \textbf{Gestión de errores}: medida en la que la aplicación presenta una baja tasa de errores y proporciona mecanismos adecuados para su resolución por parte del usuario.
    \item \textbf{Facilidad de aprendizaje}: grado de facilidad con el que un usuario puede aprender a utilizar la aplicación, ya sea mediante su intuitividad o a través de manuales de uso.
    \item \textbf{Memorabilidad}: capacidad de la aplicación para ser recordada por el usuario tras un periodo sin uso.
    \item \textbf{Satisfacción}: nivel de satisfacción general del usuario, considerando tanto la resolución de problemas (como la gestión de alertas y casos) como la experiencia estética.
\end{itemize}

La evaluación se llevará a cabo mediante una encuesta basada en la escala de Likert, en la que cada ítem se valora del 1 al 5, donde 1 indica ``totalmente en desacuerdo'' y 5 ``totalmente de acuerdo''.

Los ítems del cuestionario son los siguientes:

\begin{enumerate}
    \item AlertRes es una aplicación fácil de aprender y de utilizar.
    \item AlertRes presenta un diseño atractivo y fácil de comprender.
    \item AlertRes ha resultado ser una herramienta útil para la gestión de casos de personas desaparecidas.
    \item AlertRes permite navegar de forma rápida y sencilla a través de sus diferentes funcionalidades.
    \item Me he sentido en control en todo momento al utilizar la aplicación.
    \item Las características de la aplicación facilitan recordar su uso tras periodos prolongados sin utilizarla.
    \item El manual de instrucciones y los canales de contacto han sido útiles ante la aparición de errores o dudas.
    \item En general, estoy satisfecho/a con la aplicación.
    \item No han surgido problemas relevantes durante el uso de la aplicación.
\end{enumerate}

Finalmente, se incluye una pregunta abierta para recoger sugerencias de mejora:

\begin{itemize}
    \item Aportaciones o sugerencias de mejora:
\end{itemize}
